% ==================================================================
% id:       rv-005
% topic:    משתנים מקריים
% subtopic: משתנים מקריים בלתי תלויים
% solution: full
% verified: no
% ==================================================================
זורקים 2 מטבעות הוגנים ובלתי תלויים, שבצדדיהם $0$ ו-$1$. נגדיר את המשתנים המקריים:
$X_0$ = מספר האפסים שיצאו; \ $X_1$ = מספר האחדות שיצאו; \ $Y$ = סכום המספרים שיצאו; \ $Z=1$ אם תוצאות המטבעות שונות ו-$Z=0$ אם הן זהות.

בכל סעיף בדקו, בעזרת ההגדרה, אם שני המשתנים בלתי תלויים:
\begin{solbox}
\unverified
בכל הסעיפים המשתנים \textbf{תלויים}. בכל אחד מספיק למצוא זוג ערכים $a,b$ שעבורו
\[
\PP(U=a,\ V=b)\neq\PP(U=a)\,\PP(V=b).
\]
\end{solbox}
\begin{enumerate}
\item $X_0,X_1$;
\begin{solbox}
\unverified
$\PP(X_0=2,X_1=2)=0$ (מאורע ריק), אבל $\PP(X_0=2)\PP(X_1=2)=\frac14\cdot\frac14\neq0$.
\end{solbox}
\item $X_0,Y$;
\begin{solbox}
\unverified
$\PP(X_0=0,Y=0)=0$ (בשני המטבעות יצא 1, אבל הסכום 0), אבל $\PP(X_0=0)\PP(Y=0)=\frac14\cdot\frac14\neq0$.
\end{solbox}
\item $X_1,Z$;
\begin{solbox}
\unverified
$\PP(X_1=2,Z=1)=0$, אבל $\PP(X_1=2)\PP(Z=1)=\frac14\cdot\frac12\neq0$.
\end{solbox}
\item $Y,Z$.
\begin{solbox}
\unverified
$\PP(Y=0,Z=1)=0$, אבל $\PP(Y=0)\PP(Z=1)=\frac14\cdot\frac12\neq0$.
\end{solbox}
\end{enumerate}
